% !TEX root = ../../../main.tex
% Block card: detection.drop_detection -- "Drop detection" (Part I, the blocks chapter).
\cardfield{What it does}
Finds each sudden drop in the signal, marks where it starts and where it bottoms out, and
draws a frame around it that holds that one drop and nothing of its neighbours.

\cardfield{When to use it}
When the events you care about are steep falls: the ``sharkfins'' (a slow climb, then a
sudden fall) and ``troughs'' (a sudden fall from a flat stretch, then a slow recovery) seen on
long fungal recordings. It is the last of three blocks in the ready-made drop detector:
\blockref{preprocessing.detrend} removes the slow drift, \blockref{detection.stage_encoding}
turns the trace into letters, and this block reads the letters and the trace to find the
drops. It needs that encoding in front of it; it refuses anything else.

\cardfield{What goes in, what comes out}
In go the letters from the encoding and the detrended trace. Out comes one stretch of time
per confirmed drop, scored by how deep the drop is in millivolts. Each drop's start, its
bottom and how it was judged are kept with it, and the block page lists the candidates it
threw away and why.

\begin{plainterms}[How it decides]
Think of walking along a mountain path and noting every place the ground falls away
steeply. The block first looks only at places the encoding already marked as a fast fall, or
just after a climb. At each one it asks: is the ground really steeper here than the usual
bumpiness of the path? Where does the fall level out? Was there a real climb before it? Is
this fall the biggest thing in its own frame? A drop that fails any question is set aside and
counted.
\end{plainterms}

\cardfield{Settings that matter}
\emph{slope\_sigma} is the main one. A drop is only seen if, somewhere along its fall, the
trace drops faster than this many times its own background jitter. Lower it (to 6, 4, even 2)
when the block misses gentle drops you can see by eye; raise it to keep only the sharpest.
It must be the same number as in the encoding block before it. \emph{min\_depth\_frac}
throws away drops shallower than this fraction of the deepest drop in the span: at 0.1, a
drop a tenth the size of the biggest one is dropped. \emph{morphology} says whether the span
is made of sharkfins or troughs; \texttt{auto} decides once for the whole span.

\cardfield{Watch out for}
When the block finds too few drops, the cause is almost always the slope gate
(\emph{slope\_sigma}), not the length of the drift removal in front of it: lower
\emph{slope\_sigma} first. The two copies of \emph{slope\_sigma} (here and in the encoding)
are not checked against each other. Because the shallow-drop filter is relative to the
deepest drop, one huge drop in a span can push real smaller ones out. Gaps in the recording
are skipped: no drop is drawn across a gap, and a drop that runs into one is left out and
counted.
